\subsection{Z Transform}

\subsubsection*{Motivation and Intuition}

The Z transform is the discrete-time counterpart of the Laplace transform. It is used for sequences rather than continuous functions. This makes it essential in digital signal processing, sampled-data systems, digital control, and difference equations.

Where the Laplace transform analyzes continuous-time systems, the Z transform analyzes discrete-time systems.

\begin{conceptbox}
    For a sequence $f_n$, the one-sided Z transform is defined by
    \[
        F(z)=Z[f_n]=\sum_{n=0}^{\infty}f_nz^{-n}.
    \]
    The inverse Z transform is written as
    \[
        f_n=Z^{-1}[F(z)].
    \]
\end{conceptbox}

Here, $z$ is a complex variable. Since the Z transform is a power series in $z^{-1}$, it has a region of convergence.

\begin{conceptbox}
    The region of convergence is the set of values of $z$ for which
    \[
        \sum_{n=0}^{\infty}f_nz^{-n}
    \]
    converges.
\end{conceptbox}

\subsubsection*{Connection to the Laplace Transform}

If a continuous signal is sampled at intervals of length $T$, then we obtain values
\[
    f(0), f(T), f(2T), \ldots.
\]

Let
\[
    f_n=f(nT).
\]

The Z transform is then used to analyze the sequence:
\[
    F(z)=\sum_{n=0}^{\infty}f_nz^{-n}.
\]

The connection with the Laplace transform is commonly expressed through
\[
    z=e^{sT}.
\]

\subsubsection*{Basic Z Transforms}

\[
    \begin{aligned}
        Z[1]     & = \frac{z}{z-1}, \qquad |z|>1, \\
        Z[a^n]   & = \frac{1}{1-az^{-1}}
        = \frac{z}{z-a}, \qquad |z|>|a|,          \\
        Z[n a^n] & = \frac{az}{(z-a)^2}.
    \end{aligned}
\]

\subsubsection*{Z Transforms of Common Sequences}

\begin{conceptbox}
    For a sequence $f_n$, $n\geq 0$, the one-sided Z transform is
    \[
        F(z)=Z[f_n]=\sum_{n=0}^{\infty}f_nz^{-n}.
    \]

    Commonly used Z transform pairs are:

    \[
        \begin{array}{c|c|c}
            \textbf{Sequence } f_n,\ n\geq 0                    & \textbf{Z Transform } F(z) & \textbf{Region of Convergence} \\ \hline
            1                                                   & \dfrac{z}{z-1}             & |z|>1                          \\[8pt]
            a^n                                                 & \dfrac{z}{z-a}             & |z|>|a|                        \\[8pt]
            na^n                                                & \dfrac{az}{(z-a)^2}        & |z|>|a|                        \\[8pt]
            n                                                   & \dfrac{z}{(z-1)^2}         & |z|>1                          \\[8pt]
            \delta_n                                            &
            1                                                   & |z|>0                                                       \\[8pt]
            \delta_{n-m}                                        &
            z^{-m}                                              & |z|>0                                                       \\[8pt]
            \sin(n\omega T)                                     &
            \dfrac{z\sin(\omega T)}{z^2-2z\cos(\omega T)+1}     & |z|>1                                                       \\[12pt]
            \cos(n\omega T)                                     &
            \dfrac{z[z-\cos(\omega T)]}{z^2-2z\cos(\omega T)+1} & |z|>1
        \end{array}
    \]
\end{conceptbox}

\begin{noteBox}
    The Z transform table is the discrete-time counterpart of the Laplace transform table. It is especially useful when solving difference equations, analyzing digital filters, and studying sampled systems.
\end{noteBox}

\subsubsection*{Key Properties}

\begin{conceptbox}
    \textbf{Linearity}
    \[
        Z[c_1f_n+c_2g_n]
        =
        c_1F(z)+c_2G(z).
    \]
\end{conceptbox}

\begin{conceptbox}
    \textbf{Multiplication by $n$}
    If
    \[
        Z[f_n]=F(z),
    \]
    then
    \[
        Z[nf_n]
        =
        -z\frac{dF(z)}{dz}.
    \]
\end{conceptbox}

\begin{conceptbox}
    \textbf{Discrete Convolution}
    If
    \[
        w_n=\sum_{k=0}^{n}f_kg_{n-k},
    \]
    then
    \[
        W(z)=F(z)G(z).
    \]
\end{conceptbox}

\subsubsection*{Worked Examples}

\begin{examplebox}
    \textbf{Example 1: Z transform of the unit sequence}

    Let
    \[
        f_n=1, \qquad n\geq 0.
    \]

    Using the definition,
    \[
        F(z)=\sum_{n=0}^{\infty}1\cdot z^{-n}.
    \]

    So,
    \[
        F(z)=\sum_{n=0}^{\infty}z^{-n}.
    \]

    This is a geometric series with common ratio $z^{-1}$:
    \[
        F(z)=1+z^{-1}+z^{-2}+z^{-3}+\cdots.
    \]

    For $|z^{-1}|<1$, or equivalently $|z|>1$, the sum is
    \[
        F(z)=\frac{1}{1-z^{-1}}.
    \]

    Multiplying numerator and denominator by $z$,
    \[
        F(z)=\frac{z}{z-1}.
    \]

    Therefore,
    \[
        \boxed{
            Z[1]=\frac{z}{z-1}, \qquad |z|>1.
        }
    \]
\end{examplebox}

\begin{examplebox}
    \textbf{Example 2: Z transform of $a^n$}

    Let
    \[
        f_n=a^n, \qquad n\geq 0.
    \]

    Using the definition,
    \[
        F(z)=\sum_{n=0}^{\infty}a^nz^{-n}.
    \]

    Rewrite the product:
    \[
        a^nz^{-n}=(az^{-1})^n.
    \]

    Therefore,
    \[
        F(z)=\sum_{n=0}^{\infty}(az^{-1})^n.
    \]

    This is a geometric series with common ratio $az^{-1}$. It converges if
    \[
        |az^{-1}|<1.
    \]

    Thus,
    \[
        |z|>|a|.
    \]

    The sum is
    \[
        F(z)=\frac{1}{1-az^{-1}}.
    \]

    Multiplying numerator and denominator by $z$,
    \[
        F(z)=\frac{z}{z-a}.
    \]

    Therefore,
    \[
        \boxed{
            Z[a^n]=\frac{1}{1-az^{-1}}=\frac{z}{z-a}, \qquad |z|>|a|.
        }
    \]
\end{examplebox}

\begin{examplebox}
    \textbf{Example 3: Z transform of $\cos(nx)$}

    Using Euler's formula,
    \[
        \cos(nx)=\frac{1}{2}\left(e^{inx}+e^{-inx}\right).
    \]

    Apply linearity:
    \[
        Z[\cos(nx)]
        =
        \frac{1}{2}Z[e^{inx}]
        +
        \frac{1}{2}Z[e^{-inx}].
    \]

    Since
    \[
        Z[a^n]=\frac{1}{1-az^{-1}},
    \]
    we use
    \[
        a=e^{ix}
    \]
    and
    \[
        a=e^{-ix}.
    \]

    Thus,
    \[
        Z[e^{inx}]
        =
        \frac{1}{1-e^{ix}z^{-1}},
    \]
    and
    \[
        Z[e^{-inx}]
        =
        \frac{1}{1-e^{-ix}z^{-1}}.
    \]

    Therefore,
    \[
        F(z)
        =
        \frac{1}{2}
        \frac{1}{1-e^{ix}z^{-1}}
        +
        \frac{1}{2}
        \frac{1}{1-e^{-ix}z^{-1}}.
    \]

    Use a common denominator:
    \[
        F(z)
        =
        \frac{1}{2}
        \frac{1-e^{-ix}z^{-1}+1-e^{ix}z^{-1}}
        {(1-e^{ix}z^{-1})(1-e^{-ix}z^{-1})}.
    \]

    Simplify the numerator:
    \[
        1-e^{-ix}z^{-1}+1-e^{ix}z^{-1}
        =
        2-(e^{ix}+e^{-ix})z^{-1}.
    \]

    Since
    \[
        e^{ix}+e^{-ix}=2\cos x,
    \]
    the numerator becomes
    \[
        2-2\cos(x)z^{-1}.
    \]

    Including the factor $\frac{1}{2}$:
    \[
        \text{numerator}=1-\cos(x)z^{-1}.
    \]

    Now simplify the denominator:
    \[
        (1-e^{ix}z^{-1})(1-e^{-ix}z^{-1})
    \]
    \[
        =
        1-(e^{ix}+e^{-ix})z^{-1}+e^{ix}e^{-ix}z^{-2}.
    \]

    Since
    \[
        e^{ix}e^{-ix}=1,
    \]
    we get
    \[
        1-2\cos(x)z^{-1}+z^{-2}.
    \]

    Therefore,
    \[
        \boxed{
            Z[\cos(nx)]
            =
            \frac{1-\cos(x)z^{-1}}
            {1-2\cos(x)z^{-1}+z^{-2}}
        }.
    \]
\end{examplebox}

\begin{examplebox}
    \textbf{Example 4: Inverse Z transform using partial fractions}

    Find the inverse Z transform of
    \[
        F(z)=\frac{2z^2}{(z+2)(z+1)^2}.
    \]

    The standard method is to expand $\dfrac{F(z)}{z}$:
    \[
        \frac{F(z)}{z}
        =
        \frac{2z}{(z+2)(z+1)^2}.
    \]

    Write
    \[
        \frac{2z}{(z+2)(z+1)^2}
        =
        \frac{A}{z+2}
        +
        \frac{B}{z+1}
        +
        \frac{C}{(z+1)^2}.
    \]

    Multiply both sides by $(z+2)(z+1)^2$:
    \[
        2z=A(z+1)^2+B(z+2)(z+1)+C(z+2).
    \]

    From partial fraction computation,
    \[
        A=-4,\qquad B=4,\qquad C=-2.
    \]

    Thus,
    \[
        \frac{F(z)}{z}
        =
        \frac{-4}{z+2}
        +
        \frac{4}{z+1}
        -
        \frac{2}{(z+1)^2}.
    \]

    Multiplying by $z$,
    \[
        F(z)
        =
        -\frac{4z}{z+2}
        +
        \frac{4z}{z+1}
        -
        \frac{2z}{(z+1)^2}.
    \]

    Using known inverse transforms:
    \[
        Z^{-1}\left[\frac{z}{z+2}\right]=(-2)^n,
    \]
    \[
        Z^{-1}\left[\frac{z}{z+1}\right]=(-1)^n,
    \]
    and
    \[
        Z^{-1}\left[\frac{z}{(z+1)^2}\right]=n(-1)^{n+1}.
    \]

    Therefore,
    \[
        f_n
        =
        -4(-2)^n
        +
        4(-1)^n
        -
        2n(-1)^{n+1}.
    \]

    So,
    \[
        \boxed{
            f_n=4(-1)^n-4(-2)^n-2n(-1)^{n+1}
        }.
    \]
\end{examplebox}

\subsubsection*{Engineering Insight}

\begin{noteBox}
    The Z transform is fundamental in digital signal processing. Difference equations, digital filters, and sampled control systems are usually analyzed using the Z transform. Stability is often determined by the location of poles in the complex $z$-plane.
\end{noteBox}

\subsubsection*{Common Mistakes}

\begin{conceptbox}
    Common mistakes in Z transforms include:
    \[
        \begin{aligned}
             & \text{1. Forgetting that the Z transform is for sequences, not continuous functions.}              \\
             & \text{2. Ignoring the region of convergence.}                                                      \\
             & \text{3. Confusing } z^{-1} \text{ notation with ordinary negative powers without interpretation.} \\
             & \text{4. Expanding } F(z) \text{ instead of } \frac{F(z)}{z} \text{ during inverse transforms.}    \\
             & \text{5. Treating discrete convolution like ordinary multiplication in time.}
        \end{aligned}
    \]
\end{conceptbox}